\documentclass{article}
\usepackage{amsmath,amssymb,amsthm}
\begin{document}

\textbf{Subproblem 2 --- Opposite-sign endpoints and boundary-chain decomposition.}

\medskip
\textbf{Context.}
Use the same setup as Subproblem~1.  Assume the \emph{equality case}:
$|\{p:F(p)\neq 0\}|=n$.  By Subproblem~1 (whose results may be freely used),
the support of $F$ is exactly $\{\ell_s:s\in S\}$, all $n$ points are lonely,
and $F(\ell_s)=\varepsilon(t_s)\in\{\pm1\}$.

Assume additionally:
\begin{itemize}
  \item $|S|\ge 3$,
  \item $T$ contains at least one vector with $\varepsilon(t)=+1$ and at least one
        with $\varepsilon(t)=-1$.
\end{itemize}

\medskip
\textbf{Statement.}  Under the above hypotheses prove:
\begin{enumerate}
  \item The signs of the lonely points are not all the same: there exist
        $s_a, s_b\in S$ with $F(\ell_{s_a})=+1$ and $F(\ell_{s_b})=-1$.
  \item After an affine change of coordinates (without loss of generality) one
        may place $s_a$ at the leftmost and $s_b$ at the rightmost vertex of
        $\operatorname{conv}(S)$.  The boundary of $\operatorname{conv}(S)$ then
        decomposes into an \emph{upper chain} $U$ and a \emph{lower chain} $L$,
        each a strictly $x$-monotone convex arc from $s_a$ to $s_b$, with
        $|U|+|L|=n+2$.
  \item The sign assignment $s\mapsto \varepsilon(t_s)$ is not constant on
        either chain: each chain contains at least one vertex with sign $+1$
        and at least one with sign $-1$.
\end{enumerate}

\medskip
\textbf{Hint.}
For~(1): if all signs were equal, say all $+1$, then every translate used
would have sign $+1$ (since $\varepsilon(t_s)=+1$ for all $s$, and $t_s$ is
the only translate covering $\ell_s$).  But then every point in $S+T$ has
$F(p)\ge 0$, with equality only at uncovered points; in particular there are
no $-1$ contributions, contradicting the assumption that $\varepsilon$ attains
the value $-1$.

For~(2): choose the affine direction $x$ to be generic (not parallel to any
chord of $S$); then the leftmost and rightmost vertices are unique.

For~(3): $s_a$ has sign $+1$ (with $s_a$ itself on both chains as an endpoint)
and $s_b$ has sign $-1$, so the sign changes along each chain.

\end{document}
